Antes de suponer que los grupos de núcleos se parecen a sockets, lo medimos. Un pequeño programa lanza una pelota entre cada par de núcleos y cronometra cuánto tarda en volver.

A la izquierda, el mapa: más oscuro es más lento. Los cuadros claros de la diagonal son los cuatro clústeres de núcleos E: dentro de cada uno, los viajes son rápidos.

A la derecha, las medianas. Dentro del clúster, 35 nanosegundos; entre clústeres, 47. La sorpresa: los núcleos P, los más rápidos calculando, son los más lejanos entre sí, con 86 nanosegundos.

En los Xeon de esa generación, salir del socket cuesta unas tres veces más. Aquí, 1,35. La miniatura tiene la forma correcta, pero el castigo es más suave.
